% TOP_PROBLEM: {"id": 1241, "title": "Hall's Conjecture", "category_id": 1, "category": {"id": 1, "name": "number_theory", "display_name": "Number Theory", "description": "Properties of integers, prime numbers, Diophantine equations.", "slug": "number-theory", "order_index": 1, "created_at": "2026-07-31T15:26:25.670Z"}, "status": "open"}
% FIRST_POSTED: 2026-09-25
% ENTRY_KIND: statement_only
% !TeX program = lualatex
% Definition and sources only; this file is not an LLM solution attempt.
\documentclass[11pt]{article}
\usepackage[margin=25mm]{geometry}
\usepackage{fontspec}
\setmainfont{Latin Modern Roman}
\usepackage{amsmath,amssymb,mathtools}
\usepackage{unicode-math}
\setmathfont{Latin Modern Math}
\usepackage{xurl,hyperref}
\hypersetup{colorlinks=true,urlcolor=blue}
\setcounter{secnumdepth}{0}
\setlength{\parindent}{0pt}
\setlength{\parskip}{5pt}
\setlength{\emergencystretch}{3em}
\begin{document}
\section*{\texorpdfstring{Hall's Conjecture}{Hall's Conjecture}}\label{1241}
\subsection{Definitions and mathematical statement}
For positive integers $x,y$, exclude exact equality $y^2=x^3$. Hall's selected small-power-loss assertion is
\[
 \forall{\varepsilon}>0\ \exists c_{\varepsilon}>0\ \forall x,y\in{\mathbb{N}}:\quad
 y^2\ne x^3\Longrightarrow
 |y^2-x^3|>c_{\varepsilon} x^{1/2-{\varepsilon}}.
\]
The constant depends on ${\varepsilon}$ only, not on the particular square and cube. The absolute value allows either one to be larger. This formulation retains the arbitrary ${\varepsilon}$-loss; the claim with exponent exactly $1/2$ and one positive uniform constant is stronger and is not the displayed target. A finite collection of unusually close pairs cannot disprove a statement permitting an unspecified positive constant.
\par\smallskip\textit{Formulation and concept references:} \href{https://doi.org/10.1007/BF01140190}{[S1]}.

\subsection{Short English statement}
Apart from equality, must an integer square stay at least a constant times a nearly square-root power of x away from the cube of x?
\subsection{Sources}
\begin{itemize}
\item[S1]M. Hall, The Diophantine equation x\textasciicircum{}3-y\textasciicircum{}2=k (1971).
\url{https://doi.org/10.1007/BF01140190}.
\textit{Formulation source.}
\end{itemize}
\end{document}
